\documentclass[11pt]{article}
\input{../calcnotes.sty}
\fancyhead[C]{1.4 Parametric Equations}

\begin{document}

\noindent And now for something completely different: parametric equations. Why? The best way to understand is to look at something in motion:

\vspace{1cm}

\noindent We can describe motion in a coordinate plane, but we cannot also show the time of motion. To do that, we introduce a third variable, $t$, called a parameter. Bad news: this is challenging. Good news: we can use our calculator.

\vspace{1cm}

\noindent Describe the graph of the relation determined by $x = \sqrt{t},\ y = t,\ t \ge 0$. Indicate the direction in which the curve is traced and find the Cartesian equation.

\emptygraph

\noindent Notice how important domain is\dots we call the domain $I$ the parameter interval.

\vspace{0.5cm}

\noindent \textbf{Definition:} If $x$ and $y$ are given as functions $x = f(t),\ y = g(t)$ over an interval of $t$-values, then the set of points $(x,y) = (f(t), g(t))$ defined by these equations is a \textbf{parametric curve}. The equations are \textbf{parametric equations} for the curve.

\vspace{1.5cm}

\noindent Also notice that direction matters\dots the parameter $t$ determines where the curve starts and ends. Describe the graph of the relation determined by the following. Determine initial and terminal points and find the Cartesian equation.

\[x = 2\cos t,\quad y = 2\sin t,\quad 0 \le t \le 2\pi\]

\newpage

\noindent Graph the parametric curve: $x = 3\cos t,\quad y = 4\sin t,\quad 0 \le t \le 2\pi$

\emptygraph

\vspace{1.5cm}

\noindent At this point, you might think that parametric equations are used primarily for graphs of conics. However, we can also graph a line segment. Draw and identify the graph of the parametric curve: $x = 3t,\ y=2-2t,\ 0 \le t \le 1$

\emptygraph

\noindent Find a parameterization for the line segment with endpoints $(-2,1)$ and $(3,5)$.

\vspace{4cm}

\end{document}
